\documentclass[10pt]{article}
\usepackage[T1]{fontenc}
\usepackage[utf8]{inputenc}
\usepackage[french]{babel}
\usepackage{amsmath,amssymb,mathrsfs,mathpazo}
\usepackage{geometry}
\geometry{a4paper,margin=2cm}
\usepackage{graphicx,array,tabularx,tabularray,colortbl,multirow,multicol,booktabs,enumitem,pifont,ulem,fancybox,diagbox,cancel}
\usepackage{pstricks,pstricks-add,pst-plot,pst-tree,pst-node,pst-text,pst-func,pst-math,pst-bspline,pst-3dplot,pst-eucl,pst-solides3d}
\let\pscalculate\undefined
\usepackage{pst-bezier}
\usepackage{tikz,pgfplots,tkz-tab}
\usetikzlibrary{arrows,arrows.meta,patterns,calc,positioning,trees,decorations.markings,shapes.geometric,shapes.misc,shapes.arrows,intersections,decorations.pathmorphing,decorations.pathreplacing,backgrounds,shadows,babel}
\pgfplotsset{compat=1.18}
\usepgfplotslibrary{fillbetween}
\usepackage[np]{numprint}
\usepackage{listings,makecell,pas-tableur,twoopt}
\let\barre\undefined
\usepackage[pstricks]{tablvar}
\definecolor{darkgreen}{RGB}{0,100,0}
\providecommand{\red}{\color{red}}
\providecommand{\blue}{\color{blue}}
\providecommand{\green}{\color{green}}
\providecommand{\black}{\color{black}}
\providecommand{\white}{\color{white}}
\providecommand{\gray}{\color{gray}}
\providecommand{\cyan}{\color{cyan}}
\providecommand{\magenta}{\color{magenta}}
\providecommand{\yellow}{\color{yellow}}
\providecommand{\violet}{\color{violet}}
\providecommand{\orange}{\color{orange}}
\providecommand{\euro}{\texteuro}
\providecommand{\EUR}{\texteuro}
\providecommand{\ieme}{\textsuperscript{e}}
\providecommand{\degres}{\ensuremath{^\circ}}
\providecommand{\degree}{\ensuremath{^\circ}}
\providecommand{\vv}[1]{\overrightarrow{#1}}
\providecommand{\Coord}[1]{#1}
\providecommand{\decosix}{}
\providecommand{\starredbullet}{\ensuremath{\bullet}}
\pagestyle{empty}
\setlength{\parindent}{0pt}
\usepackage{adjustbox}\DeclareUnicodeCharacter{2236}{:}\DeclareUnicodeCharacter{2190}{\ensuremath{\leftarrow}}
\providecommand{\R}{\mathbb{R}}
\providecommand{\N}{\mathbb{N}}
\providecommand{\Z}{\mathbb{Z}}
\providecommand{\Q}{\mathbb{Q}}
\providecommand{\D}{\mathbb{D}}
\providecommand{\C}{\mathbb{C}}
\providecommand{\e}{{\mathrm{e}}}
\providecommand{\vect}[1]{\overrightarrow{#1}}
\providecommand{\vectt}[1]{\overrightarrow{\mathrm{#1}}}
\providecommand{\vvt}[1]{\overrightarrow{\mathrm{#1}}}
\providecommand{\np}[2][]{{#2~#1}}
\providecommand{\vv}[1]{{\overrightarrow{#1}}}
\renewcommand{\d}{{\mathrm{d}}}
\providecommand{\degres}{{^\circ}}
\providecommand{\degre}{{^\circ}}
\providecommand{\blue}{\color{blue}}
\providecommand{\red}{\color{red}}
\providecommand{\enskip}{\quad}
\providecommand{\llbracket}{\lbrack\!\lbrack}
\providecommand{\rrbracket}{\rbrack\!\rbrack}
\providecommand{\euro}{{\text{€}}}
\providecommand{\no}{{\text{n°}}}
\providecommand{\interval}[2]{{[#1;#2]}}
\providecommand{\intervFO}[2]{{[#1;#2[}}
\providecommand{\dotfill}{\ldots}

\providecommand{\ds}{\displaystyle}
\providecommand{\textding}[1]{#1}
\providecommand{\Oij}{\(\mathrm{O};\vec{\imath},\vec{\jmath}\)}
\providecommand{\Oijk}{\(\mathrm{O};\vec{\imath},\vec{\jmath},\vec{k}\)}
\providecommand{\Ouv}{\(\mathrm{O};\vec{u},\vec{v}\)}
\providecommand{\intervFF}[2]{\ensuremath{[#1;#2]}}
\renewcommand{\euro}{\ifmmode\text{\texteuro}\else\texteuro\fi}
\DeclareUnicodeCharacter{00B0}{\ifmmode^{\circ}\else\textdegree\fi}
\DeclareUnicodeCharacter{20AC}{\ifmmode\text{\texteuro}\else\texteuro\fi}
% Énoncés : Éducation nationale. Transcription et corrigés : APMEP, Denis Vergès et contributeurs. Découpage et classement : CoopMaths, Sébastien Lozano, Éric Elter et contributeurs. Adaptation de la présentation : Maths Explorer.
% CC BY-SA : https://coopmaths.fr/www/about/
% https://coopmaths.fr/alea/static/eam/2026/tex/eam_specifique_2026_06_etrangers_1_cor.tex
\begin{document}
\noindent\textit{APMEP / CoopMaths. CC BY-SA. Présentation adaptée par Maths Explorer.}\par\medskip
Une compagnie aérienne mène une enquête de satisfaction auprès de ses clients. Les deux questions posées lors de cette enquête sont les suivantes :\\
 \og Avez-vous acheté votre billet d'avion en agence de voyage ou par internet ? \fg{}\\
	\og Êtes-vous satisfait des services proposés par notre compagnie ? \fg{}.

	\bigskip 
	Les résultats sont donnés dans le tableau suivant :
	
	\begin{center}
		\begin{tblr}{hlines,vlines,colspec={X[1.4,l]X[c]X[c]X[c]},cells={m},row{1}={font=\bfseries},column{1}={font=\bfseries}}
			& Billet acheté par internet & Billet acheté en agence & Total \\
			Satisfait & 720 & 180 & 900 \\
			Pas satisfait & 80 & 20 & 100 \\
			Total & 800 & 200 & 1\,000 \\
		\end{tblr}
	\end{center}
	
	On choisit au hasard un client ayant répondu aux deux questions.
	
	On considère les deux événements suivants :
	\begin{itemize}[label={}]
		\item $S$ : \og Le client est satisfait \fg{} ;
		\item $I$ : \og Le client a acheté son billet par internet \fg{}.
	\end{itemize}

	\begin{color}{blue}
	On utilise l'hypothèse d'équiprobabilité : «~chacun des clients à la même chance d'être choisi~».
	
	On note $\Omega$ l'univers (les 1000 clients). $\text{Card}(\Omega)$ désigne le nombre d'élément de $\Omega$.
	\end{color}



	\begin{enumerate}
		\item Calculer la probabilité qu'un client soit satisfait.
		
		\begin{color}{blue}
	$\text{P}(S) = \dfrac{\text{Card}(S)}{\text{Card}(\Omega)} = \dfrac{900}{\np{1000}} = \dfrac{9}{10}$.
	\end{color}

		
		\item Calculer la probabilité de l'événement $S\cap I$. \\
		Interpréter le résultat à l'aide d'une phrase.
		
		\begin{color}{blue}
	$\text{P}(S\cap I) = \dfrac{\text{Card}(S\cap I)}{\text{Card}(\Omega)} = \dfrac{720}{\np{1000}} = \dfrac{72}{100} = \dfrac{18}{25}$.
	
	Il y a 72~\% de chance de choisir un client satisfait \textbf{et} ayant acheté son billet sur Internet. 
	\end{color}
		
		\item Calculer $\text{P}_{S}(I)$.
		
		\begin{color}{blue}
On calcule la probabilité qu'on choisisse un client ayant acheté son billet  sur Internet \textbf{sachant}	 qu'il était satisfait.
		
		$\text{P}_{S}(I)= \dfrac{\text{Card}(S\cap I)}{\text{Card}(S)} = \dfrac{720}{\np{900}} = \dfrac{80}{100} = \dfrac{4}{5}$
		
		\end{color}
		
		\item Les deux événements $I$ et $S$ sont-ils indépendants ? Justifier la réponse.
		
		\begin{color}{blue}
		On calcule $\text{P}(I) = \dfrac{\text{Card}(I)}{\text{Card}(\Omega)} = \dfrac{800}{\np{1000}} = \dfrac{4}{5}$
		
		Comme $\text{P}_{S}(I)=\text{P}(I)$, on peut dire que les évènenents sont indépendants. On a bien dans ce cas $\text{P}(S\cap I) = \text{P}(S) \times \text{P}(I)$.
		\end{color}
		\item Le service marketing de la compagnie aérienne affirme que $90~\%$ des clients qui ont acheté leur billet en agence sont satisfaits. Est-ce le cas ? Justifier la réponse.
		
		\begin{color}{blue}
		On calcule la probabilité qu'un client soit satisfait sachant qu'il a acheté son billet en agence.
		
		$\text{P}_{\overline{I}}(S) = \dfrac{\text{Card}(S\cap \overline{I})}{\text{Card}(\overline{I})} = \dfrac{180}{200} = \dfrac{90}{100} = 90~\%$
		
		C'est bien le cas.
		
		\end{color}
	\end{enumerate}

\end{document}
